\clearpage
\VerdierPhysicalPageAnchor{0194}%
\begingroup\small\setlength{\parskip}{0pt}\setlength{\emergencystretch}{1.5em}
\setlength{\abovedisplayskip}{4pt plus 2pt minus 2pt}%
\setlength{\belowdisplayskip}{4pt plus 2pt minus 2pt}%

\VerdierSemanticAnchor{ch3d-2-2-6-proof-cone-cont}
（$L_0$ 是 $c(m)$ 的次数 $0$ 分量。）$c(m)$ 的上同调除次数 $-1$ 外均为零，
并在次数 $-1$ 等于 $X$。态射 $L_n\longrightarrow L_n\oplus Z_{n-1}$ 的矩阵
形如 $\begin{pmatrix}\operatorname{id}\\ *\end{pmatrix}$，因而该态射有一个收缩。
所以存在态射 $v\colon Z_{n-1}\longrightarrow L_{n-1}\oplus L'_{n-2}$，使得复形

\VerdierSemanticAnchor{ch3d-2-2-6-reduced-cone}
\[
\begin{aligned}
\cdots\longrightarrow0\longrightarrow0\longrightarrow Z_{n-1}
\xrightarrow{v}L_{n-1}\oplus L'_{n-2}\longrightarrow\cdots
\longrightarrow L_1\oplus L'_0\\
\longrightarrow L_0\longrightarrow0\longrightarrow\cdots
\end{aligned}
\]
拟同构于 $c(m)$。另一方面，由性质 $(E_4)$，
$L_1\oplus L'_0\longrightarrow L_0$ 的核属于 $M$。将此核记为 $L''_1$。复形

\VerdierSemanticAnchor{ch3d-2-2-6-resolution-x}
\[
\begin{aligned}
0\longrightarrow0\longrightarrow Z_{n-1}
\longrightarrow L_{n-1}\oplus L'_{n-2}\longrightarrow\cdots
\longrightarrow L_2\oplus L'_1\\
\longrightarrow L''_1\longrightarrow X\longrightarrow0\longrightarrow\cdots
\end{aligned}
\]
因而是无环的。我们由此构造了 $X$ 的一个长度为 $n-1$ 的分解，因此
$M\text{-}\dim(X)\leq n-1$。另一方面，显然
$M\text{-}\dim(X)\geq M\text{-}\dim(Y)-1$。

\VerdierSemanticAnchor{ch3d-2-2-6-conclusion}
这就完成了引理 (2.2.6) 与命题 (2.2.7) 的证明。

\VerdierSemanticAnchor{ch3d-2-3}
\medskip\noindent\textbf{2.3. 关于分解的若干引理。}

\VerdierSemanticAnchor{ch3d-2-3-1}
\medskip\noindent\textbf{命题 2.3.1.} \textbf{[22]}\quad
\textit{a) 设 $M$ 是 $\mathcal A$ 的一个对象集合，满足性质 $(E_1)$、$(E_2)$
与 $(E_3)$（2.2.1）。设
$Y\in\operatorname{Ob}\operatorname{comp}^{-}(\mathcal A)$ 是各分量均属于
$M^\wedge$（2.2.1）的复形，并设
$m\colon X\longrightarrow Y$ 是拟同构。则存在复形
$L^\bullet\in\operatorname{Ob}\operatorname{comp}^{-}(\mathcal A)$，其各分量均属于
$M$，并存在拟同构 $w\colon L^\bullet\longrightarrow X$。}

\noindent\textit{b) 反过来，设 $M$ 是 $\mathcal A$ 的一个对象集合，满足性质
$(E_1)$。假设对每个复形
$Y\in\operatorname{Ob}\operatorname{comp}^{-}(\mathcal A)$，其各分量均属于 $M$，以及每个拟同构
$m\colon X\longrightarrow Y$，都存在复形 $L^\bullet$，其各分量均属于 $M$，并存在拟同构
$w\colon L^\bullet\longrightarrow X$。则 $M$ 满足性质 $(E_2)$ 与 $(E_3)$。}

\VerdierSemanticAnchor{ch3d-2-3-1-proof}
下面证明 a)。首先，存在复形
$X'\in\operatorname{Ob}\operatorname{comp}^{-}(\mathcal A)$ 和拟同构
$X'\longrightarrow X$（1.2.1）。因此可归结到 $X$ 是
$\operatorname{comp}^{-}(\mathcal A)$ 的对象这一情形。由 (2.1.4) 与 (2.1.5)，存在态射列

\VerdierSemanticAnchor{ch3d-2-3-1-xn-sequence}
\[
u_n\colon X_n\longrightarrow X_{n-1},
\qquad n\in\mathbf Z,
\]
\endgroup
